\section{Positive holonomy, neutral extinction, and gauge}
\label{sec:holonomy}

The cofactor is positive on edges, but its periodic product is more rigid
than positivity alone suggests.

\begin{theorem}[Positive cycle holonomy]
\label{thm:positive-holonomy}
Every closed path \(\gamma=(n_0,\ldots,n_{r-1})\) satisfies

\begin{equation}
 Q(\gamma)=
 \prod_{j=0}^{r-1}\left(1+\frac1{n_j}\right)
 \in\{2,3,\ldots\}.
\label{eq:positive-holonomy}
\end{equation}
\end{theorem}

\begin{proof}
The edge equation gives

\[
 Q(\gamma)
 =\prod_j\frac{n_j+1}{n_{j+1}}
 =\frac{\prod_j(n_j+1)}{\prod_jn_{j+1}}.
\]

Because the path closes, \(\prod_jn_{j+1}=\prod_jn_j\), giving
\cref{eq:positive-holonomy}.  The original expression is a product of
positive integers, whereas every factor in the telescoped real expression is
strictly greater than one.  Hence \(Q(\gamma)\) is an integer at least two.
\end{proof}

\begin{corollary}[No neutral periodic lift]
\label{cor:no-neutral}
The skew graph

\[
 (n,g)\longrightarrow(d,q(n,d)g)
\]

has no periodic path.  Equivalently, the identity coefficient of every
rooted group trace vanishes.
\end{corollary}

\begin{proof}
A lifted path closes in the deck coordinate exactly when its base holonomy is
one.  This is excluded by \cref{thm:positive-holonomy}.
\end{proof}

The conclusion is deliberately stronger and less useful than a selective
neutral sector: it removes every recurrent orbit.  We postpone the operator
consequence \(\det_\Phi=1\) until \cref{sec:lift}, where its trace domain and
ordinary compactness boundary are explicit.

\begin{theorem}[Exact gauge decomposition]
\label{thm:gauge}
On every allowed edge,

\begin{equation}
 q(n,d)=\frac nd\left(1+\frac1n\right).
\label{eq:gauge}
\end{equation}

Thus the multiplicative cocycle is cohomologous to the source potential
\(a(n)=1+1/n\), and

\[
 Q(\gamma)=\prod_{n\in\gamma}a(n).
\]

For \(D_ue_n=n^ue_n\), on the algebraic core,

\begin{equation}
 D_u^{-1}L_{s,u}D_u e_n=
 \sum_{\substack{d\ge2\\d\mid n+1}}
 (nd)^{-s}\left(1+\frac1n\right)^{-u}e_d.
\label{eq:operator-gauge}
\end{equation}

If \(u\in i\mathbb R\), this is a unitary conjugacy.  If
\(\Re u\ne0\), it is only an algebraic or finite-window identity unless a
separate bounded-similarity theorem is proved.
\end{theorem}

\begin{proof}
Since \(dq=n+1\),

\[
 q=\frac{n+1}{d}=\frac nd\frac{n+1}{n}.
\]

The factor \(n/d\) is a vertex coboundary and telescopes on a closed path.
For \cref{eq:operator-gauge}, conjugation changes the edge coefficient by

\[
 d^{-u}q^{-u}n^u
 =\left(\frac n{dq}\right)^u
 =\left(1+\frac1n\right)^{-u}.
\]

Finally \(D_u\) is unitary precisely when \(|n^u|=1\) for all \(n\), which
holds for purely imaginary \(u\).  Otherwise \(D_u\) or its inverse is
unbounded.
\end{proof}

The cocycle is not a pure coboundary: its periodic products are nontrivial by
\cref{thm:positive-holonomy}.  The gauge theorem says that its periodic data
are already encoded by a positive source potential, not that they disappear.

\begin{figure}[t]
\centering
\begin{tikzpicture}[
  font=\footnotesize,
  >=Latex,
  vtx/.style={circle,draw=charcoal,thick,minimum size=7mm,inner sep=1pt,fill=white},
  panel/.style={draw=charcoal,rounded corners=2mm,minimum width=4.65cm,
    minimum height=3.55cm,fill=softgray!35},
  arrow/.style={->,thick,draw=formalblue},
  deck/.style={->,thick,dashed,draw=warningamber},
  map/.style={->,thick,dotted,draw=passgreen},
  note/.style={align=center,font=\scriptsize}
]
  \node[panel] (basepanel) at (0,0) {};
  \node[font=\bfseries,anchor=north] at ([yshift=-2mm]basepanel.north)
    {base cocycle};
  \node[vtx] (n) at (-1.25,0.45) {$n$};
  \node[vtx] (d) at (1.25,0.45) {$d$};
  \draw[arrow] (n) -- node[above] {$q=(n+1)/d$} (d);
  \draw[arrow,bend left=48] (d) to node[below] {$\gamma$ closes} (n);
  \node[note] at (0,-0.95)
    {$Q(\gamma)=\prod q$\\[-1mm]
     $=\prod(1+1/n)>1$};

  \node[panel] (liftpanel) at (5.25,0) {};
  \node[font=\bfseries,anchor=north] at ([yshift=-2mm]liftpanel.north)
    {regular deck lift};
  \node[vtx,minimum size=8mm] (ng) at (4.0,0.55) {$n,g$};
  \node[vtx,minimum size=8mm] (dqg) at (6.5,0.55) {$d,qg$};
  \draw[deck] (ng) -- node[above] {$q$} (dqg);
  \node[vtx,minimum size=9mm] (nQg) at (5.25,-0.8) {$n,Qg$};
  \draw[deck,bend left=28] (dqg) to node[right] {rest of $\gamma$} (nQg);
  \draw[deck,bend left=24] (nQg) to node[below] {$\ne(n,g)$} (ng);

  \node[panel] (chipanel) at (10.5,0) {};
  \node[font=\bfseries,anchor=north] at ([yshift=-2mm]chipanel.north)
    {character resolution};
  \node[draw=formalblue,rounded corners,fill=white,align=center,
    minimum width=3.75cm,minimum height=8mm] (fiber) at (10.5,0.62)
    {$L_{s,\chi}$\\edge phase $\chi(q)$};
  \node[draw=passgreen,rounded corners,fill=white,align=center,
    minimum width=3.75cm,minimum height=9mm] (haar) at (10.5,-0.75)
    {$\displaystyle\int\overline{\chi(m)}[-\log D_\chi]\,d\chi$\\
     extracts $Q=m$};
  \draw[map] (fiber) -- (haar);

  \draw[map] ([xshift=1mm]basepanel.east) --
    node[above,font=\scriptsize] {lift} ([xshift=-1mm]liftpanel.west);
  \draw[map] ([xshift=1mm]liftpanel.east) --
    node[above,font=\scriptsize] {Fourier} ([xshift=-1mm]chipanel.west);

  \node[draw=stopred,rounded corners,align=center,font=\scriptsize,
    minimum width=4.35cm,minimum height=8mm] at (0,-2.45)
    {neutral recurrence: none\\$\Phi(\mathbb L_s^r)=0$};
  \node[draw=warningamber,rounded corners,align=center,font=\scriptsize,
    minimum width=4.35cm,minimum height=8mm] at (5.25,-2.45)
    {ordinary lift: noncompact\\infinite deck multiplicity};
  \node[draw=passgreen,rounded corners,align=center,font=\scriptsize,
    minimum width=4.35cm,minimum height=8mm] at (10.5,-2.45)
    {semifinite trace: $L^1$ iff $\Re s>1/2$\\
     scalar fibers: honest Fredholm};
\end{tikzpicture}


\caption{The factor witness on a successor--divisor edge supplies the voltage
\(q=(n+1)/d\).  A closed base path accumulates
\(Q=\prod q=\prod(1+1/n)>1\), so it never closes in the regular deck
coordinate.  Character fibers retain every base cycle and permit exact
Fourier extraction of non-neutral connected coefficients.  Ordinary
noncompactness of the lift and semifinite \(L^1\)-integrability are separate
operator statements.}
\label{fig:extension-anatomy}
\end{figure}

\Cref{fig:extension-anatomy} summarizes the three levels used below.  The
base cocycle supplies the periodic label.  The regular lift tests the neutral
class.  Character fibers retain all labels as phases and allow a chosen
non-neutral coefficient to be reconstructed by Haar orthogonality.
